\section{CPU}

% - struktura CPU
% - pouziti uint8_t, uint16_t, ..., pro reprezentaci registru (platí pro celou architekturu)
% - typ CycleCounter - jeden z rysu cross-platformu
% - CPU je implementován jako interpret
% - ukazka implementace instrukci
% - jakym zpusobem jsou implementovany adresovaci rezimy
% - pocitani cyklu

Pro reprezentaci registrů a adresaci pamětí se používají datové typy s fixní
délkou~--~\texttt{uint8\_t}, \texttt{uint16\_t}, \dots, které jsou dostupné v hlavičkovém souboru
\texttt{stdint.h}. Nejen to pomáhá šetřit paměť, ale také přispívá k dokumentaci kódu.

Pro počítání cyklů se používá beznamínkový celočíselný datový typ \texttt{CycleCounter}, velikost
kterého závisí na cílové architektuře -- v případě RP2350 je to \texttt{uint32\_t}, zatímco
desktopový backend (architektura x86\_64) používá \texttt{uint64\_t}. Velikostí těchto typů
odpovídají velikosti slova zmíněných architektur.

\subsection{Vykonávací středisko CPU}

CPU byl implementován jako interpret, jelikož procesor v NESu je poměrně jednoduchý a běží na
taktovací frekvenci cca 1,7 MHz. Interpret je realizován na bázi podmínečné konstrukce
\texttt{switch} (výpis \ref{lst:cpu_step}).

\begin{listing}[htbp]
\begin{minted}{c}
void cpu_step(Cpu *self) {
    if (self->nmi_pending) {
        cpu_handle_nmi(self);
        return;
    }
    /* Fetch a decode fáze */
    uint8_t opcode = read(self, self->PC++);
    uint16_t op_addr = decode(self, mode_table[opcode]);
    /* Execute fáze (interpret) */
    switch (opcode) {
    case OPCODE_AND_IMM:
        uint8_t data = read(self, op_addr);
        self->A &= data;
        SETZ(self->A == 0);
        SETN(self->A & 0x80);
        break;
    /* Zbytek instrukcí byl zkrácen... */
    }
}
\end{minted}
\caption{Krokovací funkce CPU.}
\label{lst:cpu_step}
\end{listing}

Pro zjištění adresovacího režimu dané instrukce při dekódování, se používá LUT tabulka adresovacích
režimů, kde indexem je operační kód instrukce. Následně funkce \texttt{decode()} na základě adresovacího
režimu zjistí adresu operandu.

Výpočet adresy operandu se řídí podle tabulky, popsané v \cite{6502-datasheet}, včetně počítání
cyklů. Prvotní návrh emulátoru CPU obsahoval tabulku, kde pro každou instrukci byl přiřazen výsledný
počet cyklů. Avšak při implementaci catch-up mechanismu pro PPU, to se stalo problémem: uprostřed
dekódování CPU může zasahovat do registru PPU, a ten musí dohnat stav CPU v tento okamžik, ale když
bude synchronizován podle cyklů procesoru před vykonáním instrukce, PPU bude opožděn vůči CPU o
několik cyklů. Když PPU bude dohánět výsledný počet cyklu aktuální instrukce, ale instrukce ještě
nebyla vykonaná, PPU bude napřed vůči PPU.
